% ed.: Consolidated from the six Lagrange--Rvachev loop reports listed in
% Appendix~\ref{sec:provenance-poly}.  All notation is normalized to the
% moment/deconvolution pair of the parent volume.

% LOCAL-NON-FOURIER-HAT: \AugmentedEvaluationMatrix denotes the endpoint-basis evaluation matrix augmented by one endpoint atom
\newcommand{\AugmentedEvaluationMatrix}{\widehat U}
% LOCAL-NON-FOURIER-HAT: \AugmentedCoefficientMatrix denotes the endpoint-basis coefficient matrix augmented by one endpoint atom
\newcommand{\AugmentedCoefficientMatrix}{\widehat B}
% LOCAL-NON-FOURIER-HAT: \OmittedMatrixIndex denotes an omitted row or column index in a maximal matrix minor
\newcommand{\OmittedMatrixIndex}[1]{\widehat{#1}}
% LOCAL-NON-FOURIER-HAT: \LiftedCoefficientProjector denotes the coefficient-space lift of the polynomial interpolation projector
\newcommand{\LiftedCoefficientProjector}[2]{\widehat{\mathcal P}_{#1}^{(#2)}}
% LOCAL-NON-FOURIER-HAT: \LiftedCoefficientDetail denotes the detail operator obtained from consecutive lifted coefficient projectors
\newcommand{\LiftedCoefficientDetail}[2]{\widehat\Delta_{#1}^{(#2)}}

\section{Lagrange cardinal synthesis and exact interpolation loops}
\label{sec:lagrange-loops}

The preceding chapters solve polynomial synthesis independently of any
choice of interpolation nodes.  Lagrange interpolation adds a finite analysis
map to that synthesis map and thereby closes an exact loop.  This section
collects the six reports that developed that observation.  Their repeated
foundations are stated once in the parent volume; the genuinely different
conclusions are organized here as consequences of one diagram.

Fix pairwise distinct nodes
\[
 X=(x_0,\ldots,x_d)\subset[-1,1],
 \qquad
 \ell_j(x)=\prod_{r\ne j}\frac{x-x_r}{x_j-x_r},
\]
and write
\[
 \mathscr D_u:=\MM_u(D)^{-1},\qquad
 m\in\PositiveIntegers,\qquad h=m^{-1},\qquad
 K_m:=\{k\in\IntegerNumbers:|k|<2m\}.
\]
The mesh is called \emph{admissible through degree $d$} when
$\TwoAdicValuation(m)\ge d$.  The canonical choice is $m=2^d$.  Notice that
$\mathscr D_u x^r=\Gd_r(x)$ in the notation of
\eqref{eq:appell-pair}.

\subsection{Cardinal polynomials as literal translates}

\begin{theorem}[Universal Lagrange--Rvachev cardinal formula]
\label{thm:lag-cardinal}
For every admissible $m$, every $0\le j\le d$, and every
$x\in[-1,1]$,
\begin{equation}
 \boxed{
 \ell_j(x)=h\sum_{k\in K_m}
 (\mathscr D_u\ell_j)(kh)\,u(x-kh).}
 \label{eq:lag-cardinal}
\end{equation}
Consequently, for arbitrary nodal data $y=(y_0,\ldots,y_d)^T$ and
$I_Xy:=\sum_jy_j\ell_j$,
\begin{equation}
 \boxed{
 (I_Xy)(x)=h\sum_{k\in K_m}
 [\mathscr D_u(I_Xy)](kh)\,u(x-kh).}
 \label{eq:lag-interpolant}
\end{equation}
Both sums are finite identities, not approximations.  The atoms in these
formulas are literal translates of the original radius-one function $u$;
there is no dilation.
\end{theorem}

\begin{proof}
Use the physical-scale identity \eqref{eq:interval-algorithm} with target
interval $[-1,1]$, origin $0$, spacing $h=m^{-1}$, and radius
$\rho=mh=1$.  It applies because $\deg\ell_j=d\le\TwoAdicValuation(m)$ and gives
\eqref{eq:lag-cardinal}; compact support leaves exactly the indices
$|k|<2m$.  Multiply by $y_j$, sum over the finite node set, and use linearity
of $\mathscr D_u$ to obtain \eqref{eq:lag-interpolant}.
\end{proof}

The coefficients are closed algebraic expressions in the nodes.  Put
\[
 \lambda_j=\left(\prod_{r\ne j}(x_j-x_r)\right)^{-1},
\]
and let $e_s^{(j)}$ be the elementary symmetric polynomial of degree $s$ in
the $d$ nodes with $x_j$ omitted.

\begin{corollary}[Nodes-only Bernoulli--Bell amplitudes]
\label{cor:lag-nodes-only}
For every $0\le j\le d$ and every $t\in\RealNumbers$,
\begin{align}
 (\mathscr D_u\ell_j)(t)
 &=\sum_{r=0}^{\Floor{d/2}}
   \gamma_{2r}\frac{\ell_j^{(2r)}(t)}{(2r)!}
 \label{eq:lag-derivative-amplitude}\\
 &=\lambda_j\sum_{n=0}^{d}(-1)^{d-n}e_{d-n}^{(j)}\Gd_n(t).
 \label{eq:lag-nodes-only}
\end{align}
Thus rational nodes give a deconvolved cardinal polynomial over
$\RationalNumbers$, and rational evaluation points give rational amplitudes.
No rationality is asserted for evaluation at an arbitrary irrational real.
More explicitly, the $\gamma_{2r}$ are complete Bell polynomials in the
negated Bernoulli cumulants of \eqref{eq:appell-pair}, so no implicit linear
solve is hidden in \eqref{eq:lag-cardinal}.
\end{corollary}

\begin{proof}
The reciprocal moment series is even, and it terminates on a degree-$d$
polynomial; this gives \eqref{eq:lag-derivative-amplitude}.  Expanding
\[
 \ell_j(x)=\lambda_j\sum_{n=0}^{d}
 (-1)^{d-n}e_{d-n}^{(j)}x^n
\]
and applying $\mathscr D_ux^n=\Gd_n(x)$ term by term gives
\eqref{eq:lag-nodes-only}.  The Bell--cumulant description is exactly the
coefficient identity following \eqref{eq:appell-pair}.
\end{proof}

\paragraph{Exact Lean counterpart.}
The source-only module \path{FabiusFunction.RvachevLagrangeNodesOnly}
formalizes this corollary completely, by composition rather than through one
monolithic wrapper theorem.  The derivative identity is
\path{eval_rvachevDeconvolvedPolynomial_lagrangeBasis_eq_sum_even_iterateDerivative},
under the displayed pairwise-distinct-node hypothesis; the raw omitted-node
expansion is
\path{eval_rvachevDeconvolvedPolynomial_lagrangeBasis_eq_nodalWeight_mul_sum_appell};
and \path{lagrangeRvachevDecoder_eq_nodalWeight_mul_sum_appell} gives its
lattice-sampled form.  Here \path{Lagrange.nodalWeight} is the displayed
$\lambda_j$, and \path{rvachevAppellPolynomial} is $\Gd_n$.

For rational nodes,
\path{rvachevDeconvolvedPolynomialRat},
\path{map_rvachevDeconvolvedPolynomialRat_lagrangeBasis}, and
\path{lagrangeRvachevDecoder_eq_ratCast} give respectively the rational
polynomial, its coefficientwise real cast, and rationality at $t=k/M$.
The last identity is totalized at $M=0$ in Lean; its use in synthesis still
has the separate positive, admissible-mesh hypotheses of
\Cref{thm:lag-cardinal}.  It does not imply rational values at irrational
$t$.

Finally,
\path{rvachevReciprocalMomentRat_eq_completeBellPolynomial_neg_centeredCumulant}
identifies the full reciprocal coefficient sequence $\gamma$ with complete
Bell polynomials in the negated full centered cumulants, while
\path{momentCumulant_rvachevRawMomentRat_even_eq_bernoulliMersenne} and the
existing \path{centeredRvachevFullCumulant_odd} give the even
Bernoulli--Mersenne formula and odd vanishing.  This is an equality of formal
coefficient sequences; it does not assert analytic convergence of a
reciprocal moment-generating function.  Accordingly
\Cref{cor:lag-nodes-only} is Exact/Complete in Lean.  This promotion does not
promote the global synthesis statement \Cref{thm:lag-cardinal} or any of the
still-human clauses of \Cref{thm:lag-right-inverse}.

\subsection{The exact right inverse and its projector}

Let $q_m=|K_m|=4m-1$ and define
\begin{equation}
 U_{ik}=u(x_i-kh),
 \qquad
 B_{kj}=h(\mathscr D_u\ell_j)(kh).
 \label{eq:lag-UB}
\end{equation}
Thus $U\in\RealNumbers^{(d+1)\times q_m}$ samples the atoms and
$B\in\RealNumbers^{q_m\times(d+1)}$ synthesizes the cardinals.
For a coefficient row $c\in\RealNumbers^{K_m}$ and a continuous function $f$, write
\[
 \operatorname{Synth}c:=\sum_{k\in K_m}c_k\,u(\,\cdot-kh),
 \qquad
 I_Xf:=\sum_{j=0}^{d}f(x_j)\ell_j.
\]
Thus $\operatorname{Synth}\colon\RealNumbers^{K_m}\to C[-1,1]$ and
$I_X\colon C[-1,1]\to\Pi_d$ have compatible types in the intertwining law
below.

\begin{theorem}[Nodal closure, basis independence, and determinant identity]
\label{thm:lag-right-inverse}
The matrices in \eqref{eq:lag-UB} satisfy
\begin{equation}
 \boxed{UB=I_{d+1}.}
 \label{eq:lag-UB-I}
\end{equation}
Consequently
\begin{equation}
 \mathcal P_X:=BU,\qquad
 \mathcal P_X^2=\mathcal P_X,\qquad
 \operatorname{ran}\mathcal P_X=\operatorname{ran}B,\qquad
 \ker\mathcal P_X=\ker U,
 \label{eq:lag-projector}
\end{equation}
and
\[
 \RankOperator\mathcal P_X=\TraceOperator\mathcal P_X=d+1.
\]
Its characteristic polynomial and synthesis intertwining law are
\begin{equation}
 \chi_{\mathcal P_X}(t)=t^{q_m-d-1}(t-1)^{d+1},
 \qquad
 \operatorname{Synth}\mathcal P_X=I_X\operatorname{Synth}.
 \label{eq:lag-projector-extra}
\end{equation}

More generally, let $p_0,\ldots,p_d$ be any basis of $\PiD$ and put
\begin{equation}
 V_{in}=p_n(x_i),
 \qquad
 \mathbf D_{kn}=h(\mathscr D_{\RvachevUp}p_n)(kh).
 \label{eq:lag-UVD}
\end{equation}
Then
\begin{equation}
 \boxed{U\mathbf D=V,\qquad B=\mathbf D V^{-1},\qquad
 \mathbf D=BV.}
 \label{eq:lag-basis-factorization}
\end{equation}
For $0\le r\le d$ one also has
\begin{equation}
 B(x_0^r,\ldots,x_d^r)^T
 =h\bigl(\Gd_r(kh)\bigr)_{k\in K_m},
 \label{eq:lag-moment-transmutation}
\end{equation}
and Cauchy--Binet gives the exact partition of one
\begin{equation}
 1=\sum_{\substack{S\subset K_m\\|S|=d+1}}
 \det U_{:,S}\det B_{S,:}.
 \label{eq:lag-cauchy-binet}
\end{equation}
\end{theorem}

\begin{proof}
Evaluate \eqref{eq:lag-cardinal} at $x=x_i$ to obtain
$(UB)_{ij}=\ell_j(x_i)=\delta_{ij}$.  Hence
$\mathcal P_X^2=B(UB)U=BU$; the range, kernel, rank, and trace assertions are
the standard consequences of a right inverse.

The $n$th column of $U\mathbf D$ consists of the nodal values of the exact synthesis
of $p_n$, hence equals the $n$th column of $V$.  Since the columns of
$V^{-1}$ are the coordinates of the cardinal polynomials in the basis
$(p_n)$, $B=\mathbf D V^{-1}$, and multiplication by $V$ gives
$\mathbf D=BV$.
Taking $p_r(x)=x^r$ proves \eqref{eq:lag-moment-transmutation}.  An
idempotent is diagonalizable with eigenvalues zero and one, so its rank gives
the characteristic polynomial.  Moreover
$\operatorname{Synth}BU=I_X\operatorname{Synth}$ because $BU$ first samples
the synthesized function and then synthesizes its cardinal interpolant.
Finally apply Cauchy--Binet to $\det(UB)=1$.
\end{proof}

\paragraph{Lean coverage boundary.}
The finite index block \path{rvachevAtomIndexSet} and its typed index
\path{RvachevAtomIndex}, together with
\path{lagrangeRvachevEncoderMatrix} and
\path{lagrangeRvachevDecoderMatrix} in
\path{FabiusFunction.LagrangeRvachevMatrix}, identify the Lean matrices as
$A=hU$ and $D=h^{-1}B$.  Hence
\path{lagrangeRvachevEncoderMatrix_mul_decoderMatrix} is an exact formal proof
of the boxed clause $UB=I_{d+1}$.  It deliberately makes no claim about
$BU$, so the projector range/kernel, rank, trace, characteristic polynomial,
intertwining, basis-factorization, moment-transmutation, and Cauchy--Binet
clauses above remain human proofs.  Accordingly
\Cref{thm:lag-right-inverse} remains only partially Lean-covered.

The factorization has an exact probabilistic normalization.  Put
\[
 \widetilde U=hU,
 \qquad
 \widetilde B=h^{-1}B.
\]

\begin{proposition}[Positive Markov encoder and forced signed decoder]
\label{prop:lag-markov}
The matrix $\widetilde U$ is nonnegative and row-stochastic,
$\widetilde B\bm1=\bm1$, and
\begin{equation}
 \widetilde U\widetilde B=I_{d+1}.
 \label{eq:lag-markov-inverse}
\end{equation}
If two distinct rows of $\widetilde U$ overlap in a strictly positive column,
then $\widetilde B$ necessarily has a negative entry.
\end{proposition}

\begin{proof}
Nonnegativity is $u\ge0$, and the partition of unity at admissible mesh gives
$h\sum_ku(x_i-kh)=1$.  Since $\sum_j\ell_j=1$ and
$\mathscr D_u1=1$, every row of $\widetilde B$ sums to one.
Equation \eqref{eq:lag-markov-inverse} is \eqref{eq:lag-UB-I}.

Suppose instead that $\widetilde B\ge0$.  Its rows are then probability
vectors.  The equation
$e_i=\widetilde U_{i,:}\widetilde B$ expresses the extreme point $e_i$ of
the probability simplex as a convex combination of those rows.  Every row
used with positive weight must therefore equal $e_i$.  A column positive in
two different encoder rows would force the same decoder row to equal two
different extreme points, a contradiction.
\end{proof}

\paragraph{Lean counterpart.}
Under the same identifications, \path{lagrangeRvachevEncoderMatrix_nonneg},
\path{sum_lagrangeRvachevEncoderMatrix_row_eq_one},
\path{sum_lagrangeRvachevDecoderMatrix_row_eq_one}, and
\path{lagrangeRvachevEncoderMatrix_mul_decoderMatrix} prove every Markov and
right-inverse clause of \Cref{prop:lag-markov}.  The general theorem
\path{exists_neg_entry_of_rightInverse_of_row_overlap} proves the finite
matrix obstruction, and
\path{exists_lagrangeRvachevDecoderMatrix_entry_neg_of_row_overlap} is its
exact Lagrange--Rvachev specialization.  Both retain the displayed strictly
positive overlap as a hypothesis; no unconditional overlap claim is made.

The sign obstruction has a quantitative form.  For a row-unital right
inverse $E$ of a nonnegative row-stochastic matrix $A$, write
\[
 V(E)=\max_k\sum_j|E_{kj}|,\qquad
 \TotalVariationDistance(a,b)=\tfrac12\|a-b\|_1.
\]

\begin{theorem}[Decoder variation and the Lebesgue barrier]
\label{thm:lag-variation}
If $AE=I$ and $E\bm1=\bm1$, then for distinct rows $i,j$,
\begin{equation}
 \boxed{
 V(E)\ge\frac{2}{\|A_{i,:}-A_{j,:}\|_1}
 =\frac1{\TotalVariationDistance(A_{i,:},A_{j,:})}.}
 \label{eq:lag-TV-lower}
\end{equation}
Some row of $E$ has total negative mass at least
\begin{equation}
 \frac1{\|A_{i,:}-A_{j,:}\|_1}-\frac12.
 \label{eq:lag-negative-mass}
\end{equation}
For the Lagrange pair $(A,E)=(\widetilde U,\widetilde B)$,
\begin{equation}
 \boxed{V(\widetilde B)\ge\Lambda_X,}
 \label{eq:lag-lebesgue-barrier}
\end{equation}
where $\Lambda_X$ is the Lebesgue constant of the node set.
\end{theorem}

\begin{proof}
From $(A_{i,:}-A_{j,:})E=e_i-e_j$,
\[
 2\le\sum_k|A_{ik}-A_{jk}|\,\|E_{k,:}\|_1
 \le\|A_{i,:}-A_{j,:}\|_1V(E),
\]
which proves \eqref{eq:lag-TV-lower}.  A row summing to one has negative mass
$(\|E_{k,:}\|_1-1)/2$, giving \eqref{eq:lag-negative-mass}.

For nodal data $y$ with $\|y\|_\infty\le1$, the synthesized interpolant is
$\sum_k(B y)_ku(x-kh)$.  Since $\sum_ku(x-kh)=h^{-1}$,
\[
 \|I_Xy\|_\infty
 \le h^{-1}\|By\|_\infty
 =\|\widetilde By\|_\infty
 \le V(\widetilde B).
\]
Taking the supremum over such $y$ gives the operator norm $\Lambda_X$ and
proves \eqref{eq:lag-lebesgue-barrier}.
\end{proof}

Thus changing from cardinal coordinates to smooth compactly supported atoms
does not regularize interpolation.  It preserves the interpolating
polynomial exactly and must carry any Runge amplification into signed decoder
coordinates.

\subsection{The local shift space and its exact sequence}

Take the canonical $m=M=2^d$ and use all literal translates meeting
$[-1,1]$:
\[
 \mathcal V_d^{\rm lit}
 :=\SpanOperator\{u(\,\cdot-k/M)|_{[-1,1]}:k\in K_M\}.
\]
Let $\mathcal T$ synthesize coefficient vectors into this space and let
$\mathcal E_X$ evaluate at the $d+1$ nodes.  Finally put
\[
 \mathcal Z_X=\{f\in\mathcal V_d^{\rm lit}:f(x_j)=0\text{ for all }j\}.
\]

\begin{theorem}[Alias/sample exact sequence]
\label{thm:lag-exact-sequence}
One has
\begin{align}
 \dim\mathcal V_d^{\rm lit}&=2M+d+1,
 &\mathcal V_d^{\rm lit}\cap\RealNumbers[x]&=\PiD,
 \label{eq:lag-lit-dimension}\\
 \dim\ker\mathcal T&=2M-d-2,
 &\dim\mathcal Z_X&=2M.
 \label{eq:lag-kernel-dimensions}
\end{align}
Moreover,
\begin{equation}
 \boxed{
 \mathcal V_d^{\rm lit}=\PiD\oplus\mathcal Z_X,}
 \label{eq:lag-direct-sum}
\end{equation}
and synthesis induces the short exact sequence
\begin{equation}
 0\longrightarrow\ker\mathcal T
 \longrightarrow\ker(\mathcal E_X\mathcal T)
 \xrightarrow{\ \mathcal T\ }\mathcal Z_X
 \longrightarrow0.
 \label{eq:lag-exact-sequence}
\end{equation}
Hence the zero eigenspace of the coefficient projector has the exact split
\begin{equation}
 4M-d-2=(2M-d-2)+2M:
 \label{eq:lag-zero-split}
\end{equation}
genuine Fourier aliases plus nonzero local splines invisible at the nodes.
\end{theorem}

\begin{proof}
Apply \cref{thm:multicell} to $[-1,1]$ with physical cell width $M^{-1}$,
so the interval contains $J=2M$ cells.  It gives
$\dim\mathcal V_d^{\rm lit}=J+d+1=2M+d+1$ and the polynomial intersection
$\PiD$.  The coefficient space has dimension $4M-1$, which gives the first
kernel dimension.

Evaluation is onto because its restriction to $\PiD$ is the usual
Vandermonde isomorphism.  Rank--nullity gives $\dim\mathcal Z_X=2M$.
The intersection $\PiD\cap\mathcal Z_X$ is zero: a degree-$d$ polynomial
with $d+1$ distinct roots vanishes.  Dimensions prove
\eqref{eq:lag-direct-sum}.

If a coefficient vector is killed by $\mathcal E_X\mathcal T$, its synthesis
lies in $\mathcal Z_X$.  Conversely every member of $\mathcal Z_X$ has a
coefficient preimage, necessarily killed by the samples.  The kernel of this
restricted synthesis is exactly $\ker\mathcal T$, which proves exactness and
the dimension split.
\end{proof}

The complete $2M-d-2$-dimensional alias sector has the explicit confluent
Fourier basis of \cref{thm:alias-family}.  The complementary $2M$-dimensional
nodal-zero sector depends on $X$; it is the precise part that interpolation
removes without synthesizing the zero function.

\subsection{The minimal dictionary and its rank-one ghost}

Transport the endpoint basis $\mathcal B_d$ of
\cref{thm:endpoint-basis} from $[0,1]$ to $[-1,1]$ and denote its
$d+2$ normalized atoms by $\Phi_0,\ldots,\Phi_{d+1}$.  Let
\[
 \AugmentedEvaluationMatrix_{ir}=\Phi_r(x_i)
\]
and let $\AugmentedCoefficientMatrix$ contain the unique endpoint-basis coordinates of the
cardinals.  Then $\AugmentedEvaluationMatrix\AugmentedCoefficientMatrix=I_{d+1}$.

For $0\le r\le d+1$, delete column $r$ from $\AugmentedEvaluationMatrix$ and row $r$ from
$\AugmentedCoefficientMatrix$, and set
\begin{equation}
 v_r=(-1)^r\det\AugmentedEvaluationMatrix_{\OmittedMatrixIndex{r}},
 \qquad
 \zeta_r=(-1)^r\det\AugmentedCoefficientMatrix_{\OmittedMatrixIndex{r}}.
 \label{eq:lag-cofactor-vectors}
\end{equation}

\begin{theorem}[Rank-one complement and canonical nodal ghost]
\label{thm:lag-rank-one}
The cofactor vectors satisfy
\begin{equation}
 \AugmentedEvaluationMatrix v=0,
 \qquad \zeta^T\AugmentedCoefficientMatrix=0,
 \qquad \zeta^Tv=1,
 \label{eq:lag-cofactor-null}
\end{equation}
and
\begin{equation}
 \boxed{I_{d+2}-\AugmentedCoefficientMatrix\AugmentedEvaluationMatrix=v\zeta^T.}
 \label{eq:lag-rank-one}
\end{equation}
For $a\in\RealNumbers^{d+2}$ put
\[
 F_a=\sum_ra_r\Phi_r,
 \qquad G_X=\sum_rv_r\Phi_r.
\]
Then $G_X(x_i)=0$ for every node and
\begin{equation}
 \boxed{F_a=I_XF_a+(\zeta^Ta)G_X.}
 \label{eq:lag-functional-ghost}
\end{equation}
In particular, $F_a$ is a degree-$d$ polynomial on the interval if and only
if $\zeta^Ta=0$.
\end{theorem}

\begin{proof}
Alternating maximal minors give a right-null vector of a full-row-rank
matrix, hence $\AugmentedEvaluationMatrix v=0$; applied to $\AugmentedCoefficientMatrix^T$ they give
$\zeta^T\AugmentedCoefficientMatrix=0$.  Cauchy--Binet and
$\AugmentedEvaluationMatrix\AugmentedCoefficientMatrix=I$ yield
\[
 1=\sum_r\det\AugmentedEvaluationMatrix_{\OmittedMatrixIndex{r}}
       \det\AugmentedCoefficientMatrix_{\OmittedMatrixIndex{r}}=\zeta^Tv.
\]
The complementary idempotent $I-\AugmentedCoefficientMatrix\AugmentedEvaluationMatrix$ has rank one, range
$\ker\AugmentedEvaluationMatrix=\SpanOperator\{v\}$, and left range
$\ker\AugmentedCoefficientMatrix^T=\SpanOperator\{\zeta\}$.  The normalization
$\zeta^Tv=1$ proves \eqref{eq:lag-rank-one}.

The endpoint coefficient vector of $I_XF_a$ is
$\AugmentedCoefficientMatrix\AugmentedEvaluationMatrix a$.  Subtracting it from $a$ and applying
\eqref{eq:lag-rank-one} proves \eqref{eq:lag-functional-ghost}.  The ghost is
nonpolynomial: if it belonged to $\PiD$, its $d+1$ nodal zeros would make it
zero, contradicting $v\ne0$ and independence of the endpoint basis.  Thus the
scalar defect (ghost) coefficient $\zeta^Ta$ vanishes exactly for polynomial
synthesis.
\end{proof}

This reconciles two apparently different ``defects.''  The universal
polynomiality functional of \cref{thm:certificate} is node independent; after
normalization it is the same one-dimensional quotient coordinate as
$\zeta^T$.  The ghost $G_X$ is a node-dependent representative of that
universal nonpolynomial quotient, chosen to vanish at the interpolation
nodes.

\begin{corollary}[Defect-completed square inverse]
\label{cor:lag-square-completion}
The square matrix
\begin{equation}
 \mathcal M_X=
 \begin{pmatrix}\AugmentedEvaluationMatrix\\ \zeta^T\end{pmatrix}
 \in\RealNumbers^{(d+2)\times(d+2)}
 \label{eq:lag-square-completion}
\end{equation}
is invertible and
\begin{equation}
 \boxed{\AugmentedCoefficientMatrix=\mathcal M_X^{-1}
 \begin{pmatrix}I_{d+1}\\0\end{pmatrix}.}
 \label{eq:lag-square-inverse}
\end{equation}
In the normalized cell coordinate $t\in[0,1]$, let
\[
 \mathcal D_d(t)=t^{d+1}+\mathcal E_d(t)
\]
be the canonical defect of \cref{thm:defect-Fourier} and let $I_\Theta$ interpolate at
the affinely transported nodes $\Theta=(\theta_0,\ldots,\theta_d)$.  The
 defect-one representative of the evaluation kernel inside the transported
 $(d+2)$-dimensional endpoint-basis space is
\begin{equation}
 \boxed{\mathcal G_{d,\Theta}
 =\mathcal D_d-I_\Theta\mathcal D_d
 =\omega_\Theta+\mathcal E_d-I_\Theta\mathcal E_d,}
 \qquad
 \omega_\Theta(t)=\prod_{j=0}^{d}(t-\theta_j).
 \label{eq:lag-canonical-ghost}
\end{equation}
Thus the algebraic ghost $v$ of \cref{thm:lag-rank-one} and the Fourier defect
of \cref{sec:certificate} are the same quotient direction with two different
normalizations.
\end{corollary}

\begin{proof}
If $\mathcal M_Xa=0$, then $\zeta^Ta=0$, so $F_a$ is a polynomial of degree
at most $d$.  The equation $\AugmentedEvaluationMatrix a=0$ says that this polynomial has the
$d+1$ distinct zeros $x_j$, hence it is zero.  Independence of the endpoint
atoms gives $a=0$, proving invertibility.  The two identities
$\AugmentedEvaluationMatrix\AugmentedCoefficientMatrix=I$ and $\zeta^T\AugmentedCoefficientMatrix=0$ are exactly
\[
 \mathcal M_X\AugmentedCoefficientMatrix=\begin{pmatrix}I\\0\end{pmatrix},
\]
which proves \eqref{eq:lag-square-inverse}.

Both $\mathcal D_d$ and $I_\Theta\mathcal D_d$ belong to the normalized local
shift space.  Their difference vanishes at every $\theta_j$, and it is
nonzero because its defect coordinate is one whereas the interpolating
polynomial has defect coordinate zero.  The evaluation kernel inside the
transported $(d+2)$-dimensional endpoint-basis space is one-dimensional, so
this is its normalized generator.  The full nodal-zero sector from
\eqref{eq:lag-exact-sequence} is much larger.  Finally,
$t^{d+1}-I_\Theta t^{d+1}=\omega_\Theta$ because both sides are monic of
degree $d+1$ and have precisely the nodes $\theta_j$ as zeros.
\end{proof}

\subsection{Nested grids and exact multiresolution algebra}

Fix an integer $s\ge0$.  Let
$X^{(0)}\subset X^{(1)}\subset\cdots\subset X^{(s)}$ be nested node
sets, with $|X^{(q)}|=d_q+1$, and let $I_q$ be the corresponding
interpolation projector.  Let $C_s$ be the coefficient space of a fixed
terminal up-dictionary.  Choose once and for all a synthesis map
$T_s:C_s\to C[-1,1]$ and its canonical coefficient map
$E_s:\Pi_{d_s}\to C_s$ at the terminal degree, so
$T_sE_s=I$ on $\Pi_{d_s}$.

\begin{theorem}[Nested projector semilattice and coefficient flag]
\label{thm:lag-nested-flag}
For all $0\le q,r\le s$,
\begin{equation}
 \boxed{I_qI_r=I_rI_q=I_{\min(q,r)}.}
 \label{eq:lag-projector-semilattice}
\end{equation}
The lifted coefficient operators
\begin{equation}
 \LiftedCoefficientProjector{q}{s}:=E_sI_qT_s
 \label{eq:lag-lifted-projector}
\end{equation}
satisfy the same multiplication law and
$\RankOperator\LiftedCoefficientProjector{q}{s}=d_q+1$.  Hence the details
\begin{equation}
 \LiftedCoefficientDetail{0}{s}=\LiftedCoefficientProjector{0}{s},
 \qquad
 \LiftedCoefficientDetail{q}{s}=
 \LiftedCoefficientProjector{q}{s}-\LiftedCoefficientProjector{q-1}{s}
 \quad(q\ge1)
 \label{eq:lag-lifted-details}
\end{equation}
are pairwise-annihilating idempotents with ranks
$d_0+1$ and $d_q-d_{q-1}$, respectively.
\end{theorem}

\begin{proof}
Assume $q\le r$.  The polynomial $I_rf$ agrees with $f$ on the smaller node
set, so $I_qI_rf=I_qf$.  Conversely $I_qf\in\Pi_{d_q}\subseteq\Pi_{d_r}$,
so $I_rI_qf=I_qf$.  This proves \eqref{eq:lag-projector-semilattice}.

Since $I_rT_s$ takes values in $\Pi_{d_s}$,
\[
 \LiftedCoefficientProjector{q}{s}\LiftedCoefficientProjector{r}{s}
 =E_sI_q(T_sE_s)I_rT_s
 =E_sI_qI_rT_s.
\]
The semilattice law follows.  The range is exactly
$E_s(\Pi_{d_q})$, proving the rank.  Expanding products of the differences
with the semilattice law proves idempotence, pairwise annihilation, and the
rank differences.
\end{proof}

\begin{corollary}[Level-grouped self-series and dyadic absolute convergence]
\label{cor:lag-self-series}
For any function $f$ for which $I_qf\to f$ in a chosen function-space
topology, put $\Delta_0=I_0$ and $\Delta_q=I_q-I_{q-1}$.  For each $q$ let
$\mathcal U_qp$ denote any finite up-atom synthesis supplied by
\cref{thm:lag-cardinal} for $p\in\Pi_{d_q}$, so that
$\mathcal U_qp=p$ on $[-1,1]$.  Then
\begin{equation}
 \boxed{\sum_{q=0}^{N}\mathcal U_q(\Delta_qf)=I_Nf
 \quad\hbox{on }[-1,1]}
 \label{eq:lag-telescoping-atoms}
\end{equation}
for every $N$.  Therefore the level-grouped atom series converges exactly in the
same topology, with exactly the same error, as the original interpolation
sequence.

Let $J_n$ denote interpolation at the $n+1$ Chebyshev--Lobatto nodes.  For
every $R>0$,
\begin{equation}
 \|u-J_nu\|_\infty=\BigO_R(n^{-R}).
 \label{eq:lag-cheb-superalg}
\end{equation}
The dyadic subsequence $I_q^{\rm CL}:=J_{2^q}$ is nested.  With
$\Delta_0^{\rm CL}=I_0^{\rm CL}$ and
$\Delta_q^{\rm CL}=I_q^{\rm CL}-I_{q-1}^{\rm CL}$, choose for each $q$ a
finite up-atom synthesis
$\mathcal U_q^{\rm CL}:\Pi_{2^q}\to C[-1,1]$ supplied by
\cref{thm:lag-cardinal}, with
$\mathcal U_q^{\rm CL}p=p$ on $[-1,1]$.  Then
\begin{equation}
 \boxed{\sum_{q=0}^{\infty}\|\Delta_q^{\rm CL}u\|_\infty<\infty,
 \qquad
 u=\sum_{q=0}^{\infty}\mathcal U_q^{\rm CL}(\Delta_q^{\rm CL}u)}
 \label{eq:lag-cheb-block-summability}
\end{equation}
absolutely and uniformly on $[-1,1]$.  Despite that summability, no uniform
geometric rate in the polynomial degree follows.
\end{corollary}

\begin{proof}
Every synthesized block equals the polynomial $\Delta_qf$ on the target
interval, so the finite sum telescopes to $I_Nf$.  This proves the exact
error transfer in any topology where the expressions are defined.

For Chebyshev--Lobatto nodes, Lebesgue's lemma gives
$\|u-J_nu\|_\infty\le(1+\Lambda_n)E_n(u)$ with
$\Lambda_n=\BigO(\log n)$.  Since the zero extension of $u$ is
$C^\infty$, Jackson estimates give $E_n(u)=\BigO_r(n^{-r})$ for every
$r$; choosing $r>R+1$ proves \eqref{eq:lag-cheb-superalg}.

The Chebyshev--Lobatto grids at degrees $2^q$ are nested, and for $q\ge1$,
\[
 \|\Delta_q^{\rm CL}u\|_\infty
 \le\|u-J_{2^q}u\|_\infty+\|u-J_{2^{q-1}}u\|_\infty
 =\BigO_R(2^{-qR}).
\]
Any $R>1$ proves summability.  On the target interval every synthesized block
equals its polynomial detail, so the Weierstrass test and finite telescoping
identity give \eqref{eq:lag-cheb-block-summability}.  A uniform
geometric rate would extend $u$ holomorphically through a Bernstein ellipse,
contradicting flatness and nontrivial compact support.
\end{proof}

At each fixed level there is no second, hidden truncation error: the
differential series defining $\mathscr D_u$ terminates at the polynomial
degree; compact support makes the lattice sum finite; and cardinal
substitution is a finite node sum.  Floating-point cancellation is a
numerical phenomenon, not an additional mathematical approximation.

\subsection{Geometric nodes and the Gaussian-binomial lift}

The interpolation geometry and the dyadic alias geometry can be varied
independently.  To see this algebraically, take geometric nodes
$1,q,\ldots,q^d$ over a characteristic-zero field, with distinct powers, and
write
\begin{align}
 \Omega_{d+1}^{(q)}(x)&=\prod_{r=0}^{d}(x-q^r)
 =\sum_{s=0}^{d+1}g_sx^s,
 \label{eq:lag-q-product}\\
 g_s&=(-1)^{d+1-s}q^{\binom{d+1-s}{2}}
 \GaussianBinomial{d+1}{s}{q}.
 \label{eq:lag-q-coeff}
\end{align}
Synthetic division gives
\begin{equation}
 \frac{\Omega_{d+1}^{(q)}(x)}{x-q^j}
 =\sum_{n=0}^{d}h_{n,j}x^n,
 \qquad
 h_{n,j}=\sum_{s=n+1}^{d+1}g_sq^{j(s-n-1)},
 \label{eq:lag-q-synthetic}
\end{equation}
and
\begin{equation}
 (\Omega_{d+1}^{(q)})'(q^j)=(-1)^j
 q^{jd-j(j+1)/2}\QPochhammer{q}{q}{j}\QPochhammer{q}{q}{d-j}.
 \label{eq:lag-q-denominator}
\end{equation}
For the literal up-function realization below, specialize to
$q\in[-1,1]$ with the powers $1,q,\ldots,q^d$ distinct, and let
$m\in\PositiveIntegers$ be an admissible integer mesh.

\begin{proposition}[Gaussian-binomial/Appell coefficient matrix]
\label{prop:lag-q-lift}
For this real specialization, $k\in K_m$, and $0\le j\le d$, the coefficient of the
$j$th geometric-node cardinal is
\begin{equation}
 \boxed{
 B_{kj}^{(q)}=
 \frac{h}{(\Omega_{d+1}^{(q)})'(q^j)}
 \sum_{n=0}^{d}h_{n,j}\Gd_n(kh).}
 \label{eq:lag-q-lift}
\end{equation}
Thus a Gaussian-binomial row from node geometry is convolved, finitely, with
the Rvachev--Appell row from moment deconvolution.  The separate dyadic
annihilator $\prod_{r=1}^{d}\QInteger{2^r}{z}$ controls compression of the resulting
atom row; it is not the same $q$-product.
\end{proposition}

\begin{proof}
The cardinal polynomial is
$\Omega_{d+1}^{(q)}(x)/
 [(x-q^j)(\Omega_{d+1}^{(q)})'(q^j)]$.  Insert
\eqref{eq:lag-q-synthetic}, apply
$\mathscr D_ux^n=\Gd_n(x)$ term by term, sample at $kh$, and multiply by the
quadrature weight $h$ from \eqref{eq:lag-cardinal}.
\end{proof}

The proposition separates the two uses of $q$-algebra cleanly: geometric
nodes create Gaussian binomial coefficients, whereas the binary sinc product
creates the $q$-integer/Thue--Morse alias law.  Their interaction is a genuine
two-parameter finite transform, not an identification of the two products.

\subsection{Appell--Vandermonde inversion and unavoidable amplification}

The scale-ladder construction of \cref{sec:ladder} gives a second, sparse
cardinal realization.  Let $y_0,\ldots,y_d$ be distinct, let
$L_j(y)$ be their cardinal polynomials, and expand
\[
 L_j(y)=\sum_{n=0}^{d}b_{nj}G_n(y).
\]
Set
\begin{equation}
 \mathsf A_{in}=G_n(y_i),\qquad
 \mathsf B=(b_{nj})_{0\le n,j\le d},\qquad
 D_n=n!\,2^{\binom{n+1}{2}}.
 \label{eq:lag-appell-vandermonde-data}
\end{equation}

\begin{theorem}[Appell--Vandermonde inverse and volume growth]
\label{thm:lag-appell-vandermonde}
One has
\begin{equation}
 \boxed{\mathsf B=\mathsf A^{-1},\qquad
 \det\mathsf A=\prod_{0\le i<j\le d}(y_j-y_i).}
 \label{eq:lag-appell-vandermonde}
\end{equation}
The raw two-atom amplitude map
\[
 \mathsf T=
 \DiagonalOperator(D_0,\ldots,D_d)\mathsf A^{-1}
\]
therefore satisfies
\begin{equation}
 \boxed{\det\mathsf T=
 \frac{2^{\binom{d+2}{3}}\prod_{n=0}^{d}n!}
 {\prod_{i<j}(y_j-y_i)}.}
 \label{eq:lag-amplitude-determinant}
\end{equation}
If the nodes lie in an interval of length two, then
\begin{equation}
 \boxed{\|\mathsf T\|_2\ge
 2^{d(d-1)/6}
 \left(\prod_{n=0}^{d}n!\right)^{1/(d+1)}.}
 \label{eq:lag-node-free-amplification}
\end{equation}
Thus even optimal node placement cannot remove the quadratic-exponential
volume expansion caused by the dyadic integration ladder.
\end{theorem}

\begin{proof}
Evaluation of the cardinal expansion at $y_i$ gives
$\mathsf A\mathsf B=I$.  Each $G_n$ is monic of degree $n$, so replacing the
monomial columns $(1,y,\ldots,y^d)$ by $(G_0,\ldots,G_d)$ is a unit-triangular
column operation.  It preserves the ordinary Vandermonde determinant and
proves \eqref{eq:lag-appell-vandermonde}.

Multiply that determinant by
\[
 \prod_{n=0}^{d}D_n
 =\left(\prod_{n=0}^{d}n!\right)
 2^{\sum_{n=0}^{d}\binom{n+1}{2}}
 =\left(\prod_{n=0}^{d}n!\right)2^{\binom{d+2}{3}}
\]
to obtain \eqref{eq:lag-amplitude-determinant}.  Finally,
$\|\mathsf T\|_2^{d+1}\ge|\det\mathsf T|$ and
$\prod_{i<j}|y_j-y_i|\le2^{\binom{d+1}{2}}$.  Taking the
$(d+1)$st root and using
\[
 \binom{d+2}{3}-\binom{d+1}{2}=\binom{d+1}{3}
\]
gives \eqref{eq:lag-node-free-amplification}.
\end{proof}

\begin{proposition}[Parity-adapted four-atom Appell ladder]
\label{prop:lag-even-four-atom-ladder}
For $n,r\ge0$, define on $[-1,1]$ the causal and anticausal Appell blocks
\begin{align}
 \Phi_n^L(x)
 &:=D_n\left[
 u\!\left(-1+\frac{x+3}{2^{n+1}}\right)
 +u\!\left(-1+\frac{x+1}{2^{n+1}}\right)
 \right]=G_n(x+2),
 \label{eq:lag-left-appell-block}\\
 \Phi_n^R(x)&:=\Phi_n^L(-x)=G_n(2-x),
 \label{eq:lag-right-appell-block}\\
 \mathscr E_r^{\rm sym}(x)
 &:=\frac12\bigl(\Phi_{2r}^L(x)+\Phi_{2r}^R(x)\bigr).
 \label{eq:lag-even-appell-block}
\end{align}
Then, for every $r\ge0$, $\mathscr E_r^{\rm sym}$ is an even monic polynomial
of degree $2r$.  Consequently, for every $N\ge0$,
$\mathscr E_0^{\rm sym},\ldots,\mathscr E_N^{\rm sym}$ form a basis of the
even polynomials of degree at most
$2N$: every such polynomial $p$ has a unique expansion
\begin{equation}
 \boxed{p(x)=\sum_{r=0}^{N}\eta_r\mathscr E_r^{\rm sym}(x),
 \qquad x\in[-1,1].}
 \label{eq:lag-even-appell-expansion}
\end{equation}
Each $\mathscr E_r^{\rm sym}$ is a symmetric four-atom up-block.  For $r=0$
the repeated atoms
combine by the partition identity, so no artificial multiplicity is required.
In particular, every even Fourier--Legendre cutoff $S_N$ has this exact
parity-adapted realization.
\end{proposition}

\begin{proof}
The two-atom identity \eqref{eq:canonical-Psi}, specialized to
$[a,b]=[-1,1]$, is exactly the first equality in
\eqref{eq:lag-left-appell-block}; reflection gives the right block.  Hence
$\mathscr E_r^{\rm sym}(-x)=\mathscr E_r^{\rm sym}(x)$.  The moment Appell
polynomial $G_{2r}$ is monic, and both
$G_{2r}(x+2)$ and $G_{2r}(2-x)$ have leading term $x^{2r}$; their average is
therefore monic of degree $2r$.  Relative to
$1,x^2,\ldots,x^{2N}$, the ordered family
$\mathscr E_0^{\rm sym},\ldots,\mathscr E_N^{\rm sym}$ has a unit-triangular
coefficient matrix.  It is therefore a
basis, proving existence and uniqueness in
\eqref{eq:lag-even-appell-expansion}.  Expanding the two definitions of
$\Phi^L$ and $\Phi^R$ displays the four atoms explicitly.
\end{proof}

\subsection{Sharp interpolation rate and endpoint amplitude barrier}

\begin{theorem}[Exact Chebyshev root-rate obstruction]
\label{thm:lag-cheb-root-rate}
Let $I_d^{\rm CL}$ be interpolation at the $d+1$ Chebyshev--Lobatto nodes.
For the nonzero compactly supported $C^\infty$ function $u$,
\begin{equation}
 \boxed{\limsup_{d\to\infty}
 \|u-I_d^{\rm CL}u\|_\infty^{1/d}=1.}
 \label{eq:lag-cheb-root-rate}
\end{equation}
Consequently there are no constants $C>0$, $\rho>1$, and $d_0$ for which
\[
 \|u-I_d^{\rm CL}u\|_\infty\le C\rho^{-d}
 \qquad(d\ge d_0).
\]
This assertion concerns an eventual uniform bound; it does not rule out
faster exceptional subsequences.
\end{theorem}

\begin{proof}
The superalgebraic convergence in \eqref{eq:lag-cheb-superalg} implies the
limsup in \eqref{eq:lag-cheb-root-rate} is at most one.  If it were less than
one, choose $\rho>1$ below its reciprocal.  The interpolation errors would
then be eventually bounded by $C\rho^{-d}$, and the best degree-$d$
approximation errors would be no larger.  Bernstein's converse theorem would
extend $u$ holomorphically to a Bernstein ellipse.  Its restriction to the
real axis would be analytic through an endpoint, impossible for a nonzero
compactly supported function that is flat there.  Hence the limsup is one.
\end{proof}

Using the constant-extension convention for $\FabiusBounded$ from the parent notation, put
\[
 \delta_d=(d+2)2^{-d}=2^{-t_d},
 \qquad t_d=d-\log_2(d+2).
\]
Define
\[
 c_{d,r}=2^{d+1}-2d-3+2r,\qquad 0\le r\le d+1.
\]
In the raw endpoint basis, write
\begin{equation}
 P(x)=\sum_{r=0}^{d+1}a_r
 u\!\left(\frac{x-c_{d,r}}{2^{d+1}}\right),
 \qquad x\in[-1,1].
 \label{eq:lag-raw-endpoint}
\end{equation}
The endpoint centers are chosen so that every displayed atom argument belongs
to $[-1,-1+\delta_d]$.

\begin{theorem}[Endpoint amplitude and quadratic-log barriers]
\label{thm:lag-endpoint-amplitude}
Every endpoint representation \eqref{eq:lag-raw-endpoint} satisfies
\begin{equation}
 \boxed{\|a\|_1\ge
 \frac{\|P\|_\infty}{\FabiusBounded(\delta_d)}.}
 \label{eq:lag-endpoint-amplitude}
\end{equation}
Consequently, for any sequence $P_d\in\Pi_d$ with
$\|P_d\|_\infty\ge1$,
\begin{equation}
 \boxed{\liminf_{d\to\infty}
 \frac{\log\|a^{(d)}\|_1}{d^2}\ge\frac{\log2}{2}.}
 \label{eq:lag-quadratic-amplitude}
\end{equation}
The same leading lower bound holds after multiplying all amplitudes by
$2^d$, as in the normalized endpoint basis.
\end{theorem}

\begin{proof}
For $d\ge2$, every active argument is $-1+s$ with
$0\le s\le\delta_d\le1$, so monotonicity and $u(-1+s)=\FabiusBounded(s)$ give
\[
 0\le
 u\!\left(\frac{x-c_{d,r}}{2^{d+1}}\right)
 \le \FabiusBounded(\delta_d)
\]
for every active atom and $x\in[-1,1]$.  For $d=0,1$ one instead uses
$0\le u\le1=\FabiusBounded(\delta_d)$, because $\delta_d\ge1$ under the stated constant
extension.  Hence in every degree
$|P(x)|\le \FabiusBounded(\delta_d)\|a\|_1$; taking the supremum proves
\eqref{eq:lag-endpoint-amplitude}.  The established endpoint law
\[
 -\log \FabiusBounded(2^{-t})=\frac{\log2}{2}t^2+\LittleO(t^2)
 \qquad(t\to\infty)
\]
is equation \texttt{dyadic-quadratic-law} in
\path{Fabius_Function_and_Rvachev_Up.tex}; its normalized limit is
kernel-verified as \path{fabiusLogProfile_normalized_tendsto} in
\path{FabiusLogSquaredAsymptotic.lean}.  This law
and $t_d/d\to1$ now give \eqref{eq:lag-quadratic-amplitude}.  A factor $2^d$
adds only $d\log2=\LittleO(d^2)$ to the logarithm.
\end{proof}

\subsection{Denominator support on an equispaced grid}

Let $d\ge1$, take the canonical synthesis mesh $m=2^d$, $h=2^{-d}$, and
put $x_j=-1+2j/d$ for $0\le j\le d$.  Define
\begin{equation}
 \mathcal P_d=
 \{p\text{ prime}:p\le d+1\}
 \cup
 \bigcup_{1\le r\le\Floor{d/2}}
 \{p\text{ prime}:p\mid 2^{2r}-1\}.
 \label{eq:lag-prime-set}
\end{equation}

\begin{theorem}[Prime support of rational cardinal amplitudes]
\label{thm:lag-prime-support}
For the equispaced grid and canonical mesh above, every prime occurring in
the reduced denominator of any literal cardinal amplitude
\[
 h(\mathscr D_u\ell_j)(kh),
 \qquad 0\le j\le d,\quad k\in K_m,
\]
belongs to $\mathcal P_d$.
\end{theorem}

\begin{proof}
After clearing the common spacing, the $j$th cardinal is an integer polynomial
divided by $2^d j!(d-j)!$.  Its coefficients and derivatives therefore
introduce only primes at most $d$.  Multiplication by $h=2^{-d}$ and
evaluation at $kh\in\IntegerNumbers[1/2]$ introduce only the already listed prime $2$.

The even cumulant
\[
 \kappa_{2r}=\frac{B_{2r}}{2r(1-4^{-r})}
\]
can have denominator primes only from $r$, from $2^{2r}-1$, and from
$B_{2r}$.  By von Staudt--Clausen, a prime in the last denominator satisfies
$p-1\mid2r$, hence $p\le2r+1\le d+1$.  Finally the reciprocal moments obey the
explicit recurrence
\begin{equation}
 \gamma_0=1,\qquad
 \gamma_n=-\sum_{r=1}^{n}
 \binom{n-1}{r-1}\kappa_r\gamma_{n-r}.
 \label{eq:lag-reciprocal-prime-recurrence}
\end{equation}
Induction shows that addition and multiplication by the displayed integers
introduce no denominator prime absent from the cumulants already used.
Substitution in \eqref{eq:lag-derivative-amplitude} proves the claim.
\end{proof}

\begin{problem}[Exact denominator valuations]
\label{prob:lag-denominator-valuations}
Determine the exact $p$-adic valuations of the least common denominator of
the whole cardinal-amplitude matrix for equispaced, dyadic, and
Chebyshev-rational nodes.  \Cref{thm:lag-prime-support} determines the possible
primes but not the exponents; the latter mix factorial valuations,
von Staudt--Clausen denominators, the factors $2^{2r}-1$, and substantial
cancellation.
\end{problem}
